% Irodov Part V §5.4, original problems 176--199 (end of polarization).
% Supplied PDF pp.225--229, printed pp.216--220; answers pp.342--343.
% Diagrams 5.32--5.35 need independent vector redraw for final edition.
\polarizationproblem{176}{216}
利用惠更斯原理，分别作出\emph{正单轴晶体}
（$n_e>n_o$）内部寻常光和非常光的波前及光线方向。
晶体光轴分别：(a) 垂直入射面并平行晶面；
(b) 位于入射面内且平行晶面；
(c) 位于入射面内，与晶面成 $45^\circ$，入射光垂直光轴。
\exsolution
对晶内每一入射波前到达晶面的点，同时以它为中心作
寻常光次波球面（速率 $v_o=c/n_o$）和非常光次波椭球面。
正晶体满足 $n_e>n_o$，所以非常波面除光轴方向外
均\emph{位于}寻常波球内，沿光轴方向两者相切。
以入射波前与界面相交处为基点，
作两个波面的公共包络切面；这两切面分别给出寻常、非常折射波前。
光线（能流方向）沿次波面相应切点的半径，不一定垂直
非常光波前。
(a) 入射面内的方向都与光轴垂直；
两波面在入射截面内可分别用速率 $c/n_o,c/n_e$
的圆表示，故有两条遵从不同折射率的折射光线。
(b) 光轴位于入射截面内，取非常光的椭圆截面作切线；
寻常光仍来自圆形截面。
(c) 非常光椭圆的长短轴相对界面倾斜；
即使波法线由界面相位匹配决定，非常光的能流方向
仍可能偏离波法线，产生走离；寻常光没有这种走离。
\excheck 所求为\emph{波面与传播方向的作图法}，
不能把非常光线一律按等效常折射率的普通斯涅耳定律处理。
本版尚须按原图恢复三幅精确作图，不能把文字算法冒充完成绘图。

\polarizationproblem{177}{216}
波长 $\lambda=589\,\mathrm{nm}$ 的自然光垂直进入由冰洲石
组成的沃拉斯顿偏振棱镜（原图5.32）。
棱镜两部分的光轴互相垂直，胶合界面的斜角
$\theta=30^\circ$。求两出射光线的夹角 $\delta$。
\exsolution
取该波长下冰洲石的近似主折射率
$n_o=1.658,\ n_e=1.486$。
在两块晶体界面上两条正交偏振分别交换所经历的
有效折射率 $n_o,n_e$；按斯涅耳定律定义
\[
 r_1=\arcsin\!\left(\frac{n_o}{n_e}\sin\theta\right),
 \qquad
 r_2=\arcsin\!\left(\frac{n_e}{n_o}\sin\theta\right).
\]
进入棱镜前端面为法向入射，内部在胶合面分别偏转
$r_1-\theta$ 和 $\theta-r_2$；
再从末端面折射到空气，两个外部偏角相加：
\[
 \boxed{\delta=
 \arcsin[n_e\sin(r_1-\theta)]
 +\arcsin[n_o\sin(\theta-r_2)]
 \simeq11^\circ}.
\]
\excheck 近似数值对冰洲石主折射率取值敏感；
实际非常光的方向与有效折射率需随原图光轴校验，
此处遵循题中两部分互换主折射率的几何模型。

\polarizationproblem{178}{216}
沿 $+z$ 传播的平面电磁波具有
$\psi=\omega t-kz$，电场的两横向分量依次为：
(a) $E_x=E\cos\psi,\ E_y=E\sin\psi$；
(b) $E_x=E\cos\psi,\ E_y=E\cos(\psi+\pi/4)$；
(c) $E_x=E\cos\psi,\ E_y=E\cos(\psi+\pi)$。
分别说明偏振态。
\exsolution
(a) 因 $E_x^2+E_y^2=E^2$，电矢量端点为圆，
故为\emph{圆偏振}；
在固定 $z$、面向来波观察时 $E$ 随时间逆时针旋转
（左右旋名称取决于观察及符号约定）。
(b) 两分量等幅、相差 $\pi/4$，消去 $\psi$ 得
\[
 \frac{E_x^2}{E^2}+\frac{E_y^2}{E^2}
 -2\cos\frac{\pi}{4}\frac{E_xE_y}{E^2}
 =\sin^2\frac{\pi}{4}.
\]
这是一条长轴沿 $y=x$ 的\emph{椭圆}，
短长半轴比为
$\tan(\pi/8)=\sqrt2-1$。
(c) $E_y=-E_x$，振动始终沿直线 $y=-x$，
为\emph{线偏振}。
\excheck 三种偏振分别为圆、椭圆、线偏振；
旋向必须连同 $z$ 轴和观察方向一起说明。

\polarizationproblem{179}{216}
从平行光轴切割的石英片中选厚度不大于
$d_{\max}=0.50\,\mathrm{mm}$ 的一片。
波长 $\lambda=589\,\mathrm{nm}$ 的入射线偏振光经过石英片，
求允许的最大片厚，使其
(a) 输出仍为线偏振（可发生偏振方向旋转）；
(b) 输出为圆偏振。
\exsolution
正入射、入射电矢量与晶体主轴成 $45^\circ$ 时，
双折射两分量相位延迟为
$\delta=2\pi|\Delta n|d/\lambda$。
取石英 $|\Delta n|\simeq0.0090$。
(a) 保持线偏振需 $\delta=m\pi$，于是
$d_m=m\lambda/(2|\Delta n|)$。
约束 $d_m\le0.50\,\mathrm{mm}$ 给最大整数 $m=15$：
\[
 \boxed{d_a=\frac{15\lambda}{2|\Delta n|}
 \simeq0.491\,\mathrm{mm}}.
\]
(b) 圆偏振需 $\delta=(2m+1)\pi/2$，
即 $d_m=(2m+1)\lambda/(4|\Delta n|)$。
允许的最大奇数为29，所以
\[
 \boxed{d_b=\frac{29\lambda}{4|\Delta n|}
 \simeq0.475\,\mathrm{mm}}.
\]
\excheck 若入射偏振方向不是与两主轴成 $45^\circ$，
只调延迟不能得到圆偏振；该条件不可略。

\polarizationproblem{180}{216}
厚度为 $d=0.50\,\mathrm{mm}$ 的石英双折射片放在
两块正交尼科耳棱镜之间，两片透振轴均与石英光轴成
$45^\circ$。设在 $0.50$--$0.60\,\mu\mathrm m$ 内
$|\Delta n|=0.0090$。求哪些波长使透射强度
不随\emph{后}检偏器绕光束旋转而改变。
\exsolution
后检偏器的透射强度与其方向无关，当且仅当
入射到它的光为圆偏振（非零强度）。
故石英片需满足四分之一波片条件
\[
 \frac{2\pi\Delta n\,d}{\lambda}
 =\frac{(2m+1)\pi}{2},\qquad
 \boxed{\lambda=\frac{4|\Delta n|d}{2m+1}}.
\]
由 $4|\Delta n|d=18\,\mu\mathrm m$，
在给定波段奇数 $2m+1$ 只能为31、33、35。
故
\[
 \boxed{\lambda=0.581,\ 0.545,\ 0.514\,\mu\mathrm m}.
\]
\excheck 对圆偏振光，无论后检偏片朝向如何，
理想片总透射其入射强度的一半。

\polarizationproblem{181}{216--217}
自然白光通过相互正交的两块尼科耳棱镜，
其中夹有沿光轴平行切割、厚度
$d=1.50\,\mathrm{mm}$ 的石英片。
晶片光轴与两透振轴都成 $45^\circ$；
在 $0.55$--$0.66\,\mu\mathrm m$ 内
$|\Delta n|=0.0090$，求透射光谱中的暗纹条数。
\exsolution
正交检偏强度
$I=\tfrac12 I_{\rm in}\sin^2(\pi\Delta n d/\lambda)$。
暗纹条件 $\Delta n d/\lambda=m$，$m$ 为正整数。
由给定波长区间，
\[
 \frac{0.0090\cdot1500}{0.66}
 \le m\le
 \frac{0.0090\cdot1500}{0.55},
\]
其中长度全以 $\mu\mathrm m$ 计；
即 $20.45\le m\le24.55$。
可取 $m=21,22,23,24$，
故
\[
 \boxed{N=4\text{ 条}}.
\]
\excheck 两端波长都不恰是整数级暗纹，无端点重复计数。

\polarizationproblem{182}{217}
平行光轴切割的晶体片厚 $d=0.25\,\mathrm{mm}$，
在 $\lambda_0=530\,\mathrm{nm}$ 时为四分之一波片。
若 $n_e-n_o=0.0090$ 近似不随可见光波长变化，
求它还在哪些可见光波长具有\emph{同号四分之一波片}作用。
\exsolution
相位延迟按整周同余后要求
$\delta\equiv\pi/2\pmod{2\pi}$（若改变手性定义则
采用 $3\pi/2$ 分支）。
对主分支条件，
\[
 \Delta n\,d=(m+\tfrac14)\lambda,\qquad
 \boxed{\lambda_m=\frac{4\Delta n d}{4m+1}}.
\]
有 $4\Delta n d=9000\,\mathrm{nm}$；
$\lambda_0\simeq9000/17=529.4\,\mathrm{nm}$ 对应 $m=4$。
可见光范围中相邻同号阶次 $m=3,5$ 给
\[
 \boxed{\lambda\simeq0.692\,\mu\mathrm m,\quad
 0.429\,\mu\mathrm m}.
\]
\excheck 相差 $3\pi/2$（反号四分之一波片）的波长
同样可实现圆偏振但手性相反；
本题沿书末的\emph{同号}波片约定列出两个答案。

\polarizationproblem{183}{217}
将石英片放在交叉尼科耳棱镜之间，晶体光轴与透振轴成
$45^\circ$。求最低片厚，使
$\lambda_1=643\,\mathrm{nm}$ 的光几乎透射最大、
$\lambda_2=564\,\mathrm{nm}$ 的光几乎完全阻断。
设两波长下 $\Delta n=0.0090$。
\exsolution
交叉偏振片的光强
$I=\tfrac12 I_0\sin^2(\pi\Delta nd/\lambda)$。
第一种光最大应有
$\Delta nd/\lambda_1=m+\tfrac12$，
第二种光最小应有
$\Delta nd/\lambda_2=k$。
由
\[
 \frac{m+1/2}{k}\simeq\frac{\lambda_2}{\lambda_1}
 =0.8771
\]
寻找最小的正整数近似，
取 $m=3,k=4$ 得 $3.5/4=0.875$。
由此
\[
 \boxed{d_{\min}\simeq\frac{3.5\lambda_1}{\Delta n}
 \simeq0.250\,\mathrm{mm}}.
\]
\excheck 此时 $\Delta nd/\lambda_2\simeq3.99$，
暗纹条件是近似满足，不能声称两色在给定有限精度下均严格极值。

\polarizationproblem{184}{217}
夹在两正交偏振片之间的石英楔，
楔角 $\theta=3.5^\circ$，晶体光轴平行楔棱，
并与偏振片透振轴成 $45^\circ$。
$\lambda=550\,\mathrm{nm}$ 的光产生间距
$\Delta x=1.0\,\mathrm{mm}$ 的干涉条纹。
求 $|\Delta n|$。
\exsolution
楔片厚度 $d(x)\simeq x\tan\theta$，
相邻同类条纹光程差增量为 $\lambda$，
故 $\Delta n\,\Delta d=\lambda$。
于是
\[
 \boxed{|\Delta n|=
 \frac{\lambda}{\Delta x\tan\theta}
 =\frac{550\times10^{-9}}
 {10^{-3}\tan3.5^\circ}
 \simeq9.0\times10^{-3}}.
\]
\excheck 小楔角下以 $\tan\theta\simeq\theta$（弧度）
替换仅带来很小误差。

\polarizationproblem{185}{217}
自然光强度 $I_0$，依次通过第一块偏振片、
引入双折射相位差 $\delta$ 的晶片与第二块偏振片。
两片透振轴分别和晶体光轴成 $\phi,\phi'$。
证明透射光强表达式，并讨论两偏振片平行和正交时的情形。
\exsolution
第一片后线偏振光强为 $I_0/2$。
在晶体寻常、非常主轴上，复振幅投影分别为
$\cos\phi$ 和 $\sin\phi$，
经晶体后相对相位为 $\delta$。
在第二片透振轴上的归一化振幅为
\[
 A=\cos\phi\cos\phi'
 +e^{i\delta}\sin\phi\sin\phi'.
\]
取模平方并利用
$\cos^2(\phi-\phi')
=(\cos\phi\cos\phi'+\sin\phi\sin\phi')^2$，得到
\[
 \boxed{I=\frac{I_0}{2}
 \left[\cos^2(\phi-\phi')
 -\sin2\phi\sin2\phi'\sin^2\frac\delta2\right]}.
\]
若平行 $\phi'=\phi$，则
\[
 \boxed{I_{\parallel}=\tfrac12 I_0
 [1-\sin^22\phi\sin^2(\delta/2)]}.
\]
若正交 $\phi'=\phi+\pi/2$，则
\[
 \boxed{I_{\perp}=\tfrac12 I_0
 \sin^22\phi\sin^2(\delta/2)}.
\]
\excheck 当 $\delta=0$ 时结果退化为
两个理想偏振片的马吕斯定律；并且
$I_{\parallel}+I_\perp=I_0/2$。

\polarizationproblem{186}{217}
圆偏振单色光垂直通过一块沿光轴平行切割的晶体薄片，
再通过尼科耳检偏棱镜。
棱镜透振轴与晶体光轴成 $\phi$；
晶片使寻常、非常分量发生相位差 $\delta$。
证明透射强度的角度依赖。
\exsolution
将右旋或左旋圆偏振电矢量的琼斯列向量
写作 $(1,\pm i)^{\mathsf T}/\sqrt2$。
晶片在两主轴上的作用为
$\operatorname{diag}(1,e^{i\delta})$；
分析器沿 $(\cos\phi,\sin\phi)$ 投影，强度为
\[
 I=\frac{I_{\rm inc}}2
 \left|\cos\phi\pm i e^{i\delta}\sin\phi\right|^2
 =\frac{I_{\rm inc}}2[1\mp\sin2\phi\sin\delta].
\]
因旋向、两晶体主轴标号可逆，
通常写为
\[
 \boxed{I=I_{\rm ref}(1+\sin2\phi\sin\delta)},
 \qquad I_{\rm ref}=I_{\rm inc}/2.
\]
\excheck 原题所要求形式的 $I_0$ 表示\emph{基准检偏强度}
时才能写成 $I_0(1+\cdots)$；
若 $I_0$ 表示晶片前总入射强度，则必须保留 $1/2$。
验证 $\delta=0$ 时 $I=I_{\rm inc}/2$。

\polarizationproblem{187}{218}
如何用理想偏振片与\emph{正单轴晶体}
（$n_e>n_o$）制成的四分之一波片，区分：
(a) 左、右旋圆偏振光；
(b) 自然光、圆偏振光以及自然光和圆偏振光的混合光？
\exsolution
先用已知快慢轴方向的四分之一波片，
把入射圆偏振态变成线偏振态，
然后旋转检偏片寻找消光方向。
两种相反旋向对应的出射线偏振方向
相对波片光轴分别为 $+45^\circ$ 和 $-45^\circ$，
因此在固定参考轴与观察方向约定后即可判旋向。
(b) 对三类光均先通过波片再旋转检偏片：
自然光始终具有与角度无关的强度；
纯圆偏振光变成完全线偏振，旋转检偏片可出现\emph{完全消光}；
自然光与圆偏振光混合后成为部分线偏振，
光强随角度变化但其最小值严格大于零。
\excheck 必须先放四分之一波片，
仅用一个线检偏片无法区分自然光与圆偏振光。

\polarizationproblem{188}{218}
两交叉偏振片 $P,A$ 中间放巴比涅补偿器 $C$
（原图5.33），它由两片光轴相互正交的石英楔组成，
小楔角 $\theta$，晶轴与检偏轴成 $45^\circ$。
插入待测双折射样品 $S$ 后，
暗纹向上移动 $\delta x$。
求 (a) 条纹间距 $\Delta x$；
(b) 样品引入的寻常与非常光程差的大小及符号。
\exsolution
取坐标 $x$ 向上为正，两石英楔厚度之差沿
$x$ 的梯度为 $2\tan\theta\simeq2\theta$。
记石英双折射量为
$\Delta n=n_e-n_o$，则补偿器光程差
$D_C(x)=2\Delta n\,\theta x$（参考常数略）。
相邻暗纹满足光程差改变量为 $\lambda$，故
\[
 \boxed{\Delta x=\frac{\lambda}{2|\Delta n|\theta}}.
\]
加入样品引入常数 $D_S$ 后，
同一暗纹条件
$D_C(x+\delta x)+D_S=D_C(x)$，
所以
\[
 \boxed{D_S=-2\Delta n\,\theta\,\delta x}.
\]
对于本约定的向上正方向，
$\Delta n>0,\delta x>0$ 意味着 $D_S<0$。
\excheck 样品光轴若与图中相反，光程差的\emph{符号}定义须反转；
原题图示需在定稿时复绘并附坐标约定。

\polarizationproblem{189}{218}
查附录，求 $\lambda=589.5\,\mathrm{nm}$ 时
石英对右、左旋圆偏振分量的折射率差。
\exsolution
设两圆偏振本征波的折射率差为
$\Delta n=n_R-n_L$。
线偏振可分解为等幅的左右旋圆偏振，
传播长度 $l$ 产生相对相移
$\Delta\varphi=(2\pi/\lambda)\Delta n\,l$。
线偏振面因此旋转其相移的一半，故
\[
 \alpha=\frac{\varphi}{l}=\frac{\pi\Delta n}{\lambda},
 \qquad \boxed{\Delta n=\frac{\alpha\lambda}{\pi}}.
\]
采用石英在该波长的旋光常数约
$\alpha=21.7^\circ/\mathrm{mm}
 =3.79\times10^2\,\mathrm{rad/m}$，
得
\[
 \boxed{|\Delta n|\simeq7.1\times10^{-5}}.
\]
\excheck 只有旋光常数转成\emph{弧度/米}
后，折射率差才无量纲。

\polarizationproblem{190}{218}
$\lambda=0.59\,\mu\mathrm m$ 的线偏振光经过
顶角 $\theta=30^\circ$ 的石英三棱镜（原图5.34），
光在石英内部沿光轴传播。
后检偏片形成条纹间隔
$\Delta x=15.0\,\mathrm{mm}$ 的明暗纹。
求石英的旋光常数及横向强度分布。
\exsolution
设 $x$ 为从某一透射极大向楔厚增加方向的距离，
厚度增加量为 $d(x)=x\tan\theta$。
石英使偏振面旋转
$\varphi(x)=\alpha d(x)$；
通过后检偏片的强度按马吕斯定律变化：
\[
 I(x)=I_{\max}\cos^2(\alpha x\tan\theta).
\]
相邻极大之距满足
$\alpha\,\Delta x\tan\theta=\pi$，所以
\[
 \boxed{\alpha=\frac{\pi}{\Delta x\tan\theta}
 \simeq3.63\times10^2\,\mathrm{rad/m}
 \simeq20.8^\circ/\mathrm{mm}},
 \quad
 \boxed{I(x)=I_{\max}\cos^2(\pi x/\Delta x)}.
\]
\excheck 旋光角的周期为 $\pi$ 而非 $2\pi$，
因为 $\cos^2$ 的周期是 $\pi$。

\polarizationproblem{191}{218}
自然单色光通过交叉尼科耳棱镜，
中间放置垂直光轴切割的石英片。
石英旋光常数 $\alpha=17^\circ/\mathrm{mm}$，
求至少多厚能使透射能流占原入射自然光的
$\eta=0.30$。
\exsolution
第一尼科耳片透射 $I_0/2$，
石英将线偏振旋转 $\alpha d$，
故
\[
 I_{\rm out}=\frac{I_0}{2}\sin^2(\alpha d).
\]
设 $I_{\rm out}/I_0=\eta=0.30$，
$\sin^2(\alpha d)=2\eta=0.60$。
取最小正根，
\[
 \boxed{d_{\min}=
 \frac{\arcsin\sqrt{2\eta}}{\alpha}
 =\frac{50.77^\circ}{17^\circ/\mathrm{mm}}
 \simeq2.99\,\mathrm{mm}}.
\]
\excheck 必须有 $\eta\le1/2$；
书末答案抽取层中的反正弦根号不可省略。

\polarizationproblem{192}{218--219}
在两个交叉尼科耳棱镜之间放一块石英片，
它在波长 $\lambda_1=436\,\mathrm{nm}$ 和
$\lambda_2=497\,\mathrm{nm}$ 下的旋光常数
分别为 $\alpha_1=41.5^\circ/\mathrm{mm}$
和 $\alpha_2=31.1^\circ/\mathrm{mm}$。
求能近似\emph{完全阻断}第一种光、
并使第二种光在自然入射时透射原光强一半的最小片厚。
\exsolution
对自然光，出射强度为
$\tfrac12 I_0\sin^2(\alpha d)$。
第一种光全暗要求
$\alpha_1d=m\,180^\circ$，
第二种光达到自然光条件下最大透射
$I_0/2$，要求
$\alpha_2d=90^\circ+k180^\circ$。
取 $m=2$ 得
\[
 \boxed{d=\frac{360^\circ}{41.5^\circ/\mathrm{mm}}
 \simeq8.67\,\mathrm{mm}}.
\]
此时 $\alpha_2d\simeq269.8^\circ$，
近于 $270^\circ$ 的最大透射条件。
\excheck 用书中两旋光常数的三位有效数只能近似同时满足
两条件；不得说精确等式恒成立。

\polarizationproblem{193}{219}
沿有轻微浑浊的圆柱玻璃管轴传播
$\lambda=589\,\mathrm{nm}$ 的线偏振光，
管内糖溶液浓度为 $500\,\mathrm{g/L}$。
侧面观察的散射亮暗螺旋条纹中，
相邻暗纹沿轴方向相隔 $l=50\,\mathrm{cm}$。
解释条纹及计算比旋光度。
\exsolution
糖溶液使线偏振面连续转动，
浑浊微粒的侧向散射强度取决于瞬时电矢量
相对观察方向的投影。旋光面每转过 $\pi$，
侧向光强分布重复一次，因此
$\alpha l=180^\circ$。
浓度换算为
$C=500\,\mathrm{g/L}=0.500\,\mathrm{g/cm^3}$，
距离 $l=5\,\mathrm{dm}$。
依比旋光度定义 $\alpha=[\alpha]C$，
\[
 \boxed{[\alpha]=\frac{180^\circ}{5\times0.500}
 =72^\circ/(\mathrm{dm}\cdot\mathrm{g/cm^3})}.
\]
\excheck 在可旋转的线偏振态下 $\pi$ 旋转回到同一偏振方向，
故间隔对应 $180^\circ$ 而非 $360^\circ$。

\polarizationproblem{194}{219}
克尔盒置于相交叉的尼科耳棱镜之间，
盒内电场与两个透振轴都成 $45^\circ$。
电容器长 $l=10\,\mathrm{cm}$，
填充硝基苯；$\lambda=0.50\,\mu\mathrm m$，
克尔常数 $B=2.2\times10^{-10}\,\mathrm{cm/V^2}$。
(a) 求透过克尔盒的光对后检偏器转角不敏感时的
最小场强；(b) 若外加 $10\,\mathrm{MHz}$
正弦电压、最大场强 $E_m=50\,\mathrm{kV/cm}$，
每秒出现多少次透射中断？
\exsolution
克尔效应引入相位差
$\delta=2\pi BlE^2$。
(a) 后检偏角不再影响透射光强，当克尔盒输出圆偏振，
最小相位差 $\delta=\pi/2$，于是
\[
 \boxed{E_{\min}=\frac1{\sqrt{4Bl}}
 \simeq10.7\,\mathrm{kV/cm}}.
\]
(b) 设 $E(t)=E_m\sin(2\pi\nu t)$，
则 $\delta(t)=\delta_m\sin^2(2\pi\nu t)$，
$\delta_m=2\pi BlE_m^2=11\pi$。
交叉检偏器的光强与 $\sin^2[\delta(t)/2]$ 成正比，
故 $\delta=2\pi k$ 时黑暗。
在一个电压周期中，$k=0$ 的两次过零
及 $k=1,\ldots,5$ 的每级四次过零，
合计 $2+4\times5=22$ 次中断，因而
\[
 \boxed{f_{\rm interrupt}=22\nu
 =2.2\times10^8\,\mathrm{s^{-1}}}.
\]
\excheck 电压平方使克尔相位具有电压半周期，
不能简单用 $2\nu$ 作答案。

\polarizationproblem{195}{219}
角频率为 $\omega$ 的线偏振光沿均匀纵向磁场
$H$ 穿过材料。若该物质费尔德常数为 $V$，
求右旋和左旋圆偏振本征波折射率之差。
\exsolution
将线偏振光分解为等幅左右圆偏振。
走过长度 $l$ 的相位差为
$\Delta\delta=(\omega/c)\Delta n\,l$；
线偏振面转角等于相位差的一半。
又法拉第旋光 $\varphi=VHl$，因而
\[
 VHl=\frac{\omega\,\Delta n\,l}{2c},
 \qquad
 \boxed{\Delta n=\frac{2cVH}{\omega}}.
\]
\excheck $\Delta n$ 符号取决于左右旋和 $H$ 方向的约定，
式子给与 $\varphi=VHl$ 一致的有向差。

\polarizationproblem{196}{219}
长度 $l=30\,\mathrm{cm}$ 的管子装有某物质，
置于纵向螺线管磁场内。
磁场强度 $H=56.5\,\mathrm{kA/m}$；
一方向场下旋光角
$\phi_1=+5^\circ10'$，反向场下
$\phi_2=-3^\circ20'$。求费尔德常数 $V$。
\exsolution
设无磁场的自然旋光角为 $\phi_{\rm nat}$，
则
$\phi_1=\phi_{\rm nat}+VHl$，
$\phi_2=\phi_{\rm nat}-VHl$。
两式相减消去自然旋光：
\[
 \boxed{V=\frac{\phi_1-\phi_2}{2Hl}}.
\]
差为 $8^\circ30'=510'$，
$2Hl=2(56500)(0.30)=33900\,\mathrm A$，
故
\[
 \boxed{V\simeq0.0150\ {\rm arcmin/A}}.
\]
\excheck 两方向旋光角的平均给自然旋光角，
不能只用 $\phi_1$ 直接除 $Hl$。

\polarizationproblem{197}{219}
原图5.35所示装置使线偏振光在右旋物质中多次往返，
物质长 $l$，比旋光系数为 $\alpha$，
费尔德常数为 $V$，纵向磁场为 $H$。
求最终透射光的偏振面总转角；
原图中光经过材料的次数为 $N=5$。
\exsolution
自然旋光属于互易效应，光沿原路反向传播时，
相对于固定外部坐标测得的旋转可相互抵消；
法拉第旋光是\emph{非互易}效应，
往返穿越的贡献可按同一实验转角方向累加。
按原图的光路和顺时针为正的约定，
自然旋光净贡献为 $\alpha l$，
磁致旋光总贡献为 $-NVHl$，所以
\[
 \boxed{\varphi=(\alpha-NVH)l,\qquad N=5}.
\]
\excheck 符号依原图的磁场箭头和观察方向；
多次反射光路不能将自然旋光也机械地乘以 $N$。

\polarizationproblem{198}{220}
装有苯的长度 $l=26\,\mathrm{cm}$ 的管子置于纵向
磁场内，夹在主方向相差 $45^\circ$ 的两块偏振片之间。
求在 $\lambda=589\,\mathrm{nm}$ 下使它作为
只沿一个方向透光的光隔离器所需的最小磁场，
并说明磁场反向的影响。
\exsolution
法拉第效应具有非互易性。
为使正向光的偏振旋转与分析器透振轴相差的
$45^\circ$ 相符，
最小正向旋转角应满足
$VHl=\pi/4$，于是
\[
 \boxed{H_{\min}=\frac{\pi}{4|V|l}}.
\]
采用附录中苯于给定波长的费尔德常数，
\[
 \boxed{H_{\min}\simeq4.0\,\mathrm{kA/m}}.
\]
磁场反向时旋光符号反向，
因此允许通过与受阻挡的传播方向互换。
\excheck 这里利用的是磁致旋光\emph{非互易性}，
不是普通互易光学旋光；若缺少两端检偏器方向，
只写 $45^\circ$ 并不足以指定正向。

\polarizationproblem{199}{220}
波长 $\lambda=0.70\,\mu\mathrm m$、
强度 $I=10\,\mathrm{W/cm^2}$ 的圆偏振光
正入射质量 $m=10\,\mathrm{mg}$ 的均匀黑色圆盘，
圆盘可绕中心法线自由转动。
圆偏振光的角动量流密度为 $I/\omega$。
经过多久圆盘角速度达到
$\omega_0=1.0\,\mathrm{rad/s}$？
\exsolution
设圆盘半径为 $R$，吸收功率为 $I\pi R^2$，
所受扭矩
$\tau=I\pi R^2/\omega$，
而圆盘转动惯量 $J=mR^2/2$。
角加速度
\[
 \dot\omega_{\rm disk}
 =\frac{\tau}{J}
 =\frac{2\pi I}{m\omega}.
\]
用光角频率 $\omega=2\pi c/\lambda$，
得
\[
 \boxed{t=\frac{m c\omega_0}{I\lambda}}.
\]
数值 $m=10^{-5}\,\mathrm{kg}$，
$I=10^5\,\mathrm{W/m^2}$，
$\lambda=7.0\times10^{-7}\,\mathrm m$，
故
\[
 \boxed{t\simeq4.29\times10^4\,\mathrm s
 \simeq11.9\,\mathrm h}.
\]
\excheck 圆盘半径严格消去，符合题目未给半径的条件；
这里忽略摩擦、反射与空气阻力。
